Full local top-$c$ retention establishes correctness, but can produce
$Gc$ intermediate survivors for only $c$ global outputs. The shared-bound
component (v63) reduces this excess selection work. The scheduler (v65)
then chooses how much scan parallelism to expose. Both operate at a fixed
refinement budget; choosing that budget is the separate problem in
\Cref{ch:refinement}.

\subsection{A complete-key upper bound with explicit witnesses}

Let $K_{(c)}$ be the global $c$th key when $|U|\geq c$. If a group holds
$c$ distinct, fully scored candidates, the largest key among those
witnesses, $\tau_g$, satisfies $K_{(c)}\preceq\tau_g$. There are already
$c$ candidates no later than $\tau_g$ in the global universe. The shared
bound is initialised to a top sentinel and updated by
\begin{equation}
 \tau\leftarrow\min_{\prec}(\tau,\tau_g).
 \label{eq:5-shared-update}
\end{equation}
No bound is published from fewer than $c$ valid candidates. If the entire
universe has fewer than $c$, no group can publish such a bound and all
valid candidates remain eligible.

\begin{proposition}[Safe rejection by a witnessed bound]
\label{prop:5-shared-bound}
Assume fixed keys, publications backed by $c$ distinct candidates, and
reads that return either the initial sentinel or one complete published
key. Rejecting a candidate only when $K_q(x)\succ\tau$ cannot remove a
global top-$c$ candidate, irrespective of publication timing.
\end{proposition}
\begin{proof}
Every published bound is at least $K_{(c)}$, and the minimum of such
bounds is also at least $K_{(c)}$. A global winner satisfies
$K_q(x)\preceq K_{(c)}\preceq\tau$, so it is never rejected. An older
bound remains an upper bound and merely permits extra work. Publications
made after earlier safe rejections remain valid: any $c$ retained
distinct candidates still supply $c$ witnesses, whether or not they are
the best $c$ of their group.
\end{proof}

This last distinction matters. With sharing enabled, a group need not
return its own full local top-$c$: it can discard locally competitive
candidates already certified as globally uncompetitive. Correctness
requires the union still to contain every global winner. The truncated
merge argument of \Cref{prop:5-local-coverage} then applies to that union.
Candidates equal to the shared bound remain eligible, including the
witness that established it. The implemented admission condition is
\begin{equation}
 K_q(x)\prec\lambda_g
 \quad\text{and}\quad
 K_q(x)\preceq\tau_{\rm seen},
 \label{eq:5-admission}
\end{equation}
where $\lambda_g$ is the local compaction threshold and
$\tau_{\rm seen}$ is the group's last observed shared bound.

\subsection{Publication, filtering and implementation limits}

For nonnegative FP32 scores generated as sums of squares, the code packs
the distance bits and nonnegative position into an unsigned 64-bit key:
\begin{equation}
 E(D,p)=2^{32}\operatorname{uintbits}_{32}(D)+p.
 \label{eq:5-packed-key}
\end{equation}
Unsigned integer comparison then agrees with the lexicographic order.
This relies on positive-zero encoding and excludes NaNs and negative
scores; it is not a general signed-float encoding. Positive infinity
orders after finite scores, and the all-one initial word orders after
valid encodings. Invalid positions are excluded from bound publication.
These are score-domain and representation preconditions, not properties
inferred from unchanged recall.

\begin{figure}[tbp]
\centering
% Algorithm diagram; the whole-key-read premise is discussed in Section 5.5.
\begin{tikzpicture}[
  font=\small\linespread{1}\selectfont,
  box/.style={draw=ThesisNavy!65,rounded corners=2pt,align=center,
    inner sep=5pt,text width=6.2cm,minimum height=.9cm},
  side/.style={box,text width=4.1cm,fill=ThesisNavy!5},
  arr/.style={-{Latex[length=2mm]},draw=ThesisNavy!80,thick}]
\node[box] (score) at (0,0) {Read primary code; compute full $\primarydist$\\
  Form $K=(\primarydist,p)$};
\node[box,fill=ThesisNavy!5] (admit) at (0,-1.8)
  {Admission: $K\prec\lambda_j$ and $K\preceq\tau_{\rm seen}$};
\node[box] (compact) at (0,-3.8) {Append admitted keys\\
  Compact when necessary\\Keep at most $c$};
\node[box] (trim) at (0,-5.9) {Final compaction and prefix trim\\
  Emit valid keys $K\preceq\tau_{\rm seen}$ to merge};
\node[side] (shared) at (6.3,-.1) {Per-query shared word $Q$\\
  Encodes bound $\tau$\\Initially all ones\\$\tau\succeq K_{(c)}$};
\node[side] (publish) at (6.3,-3.8) {Only with $c$ valid witnesses:\\
  $\operatorname{atomicMin}(Q,\kappa(\tau_j))$};
\draw[arr] (score) -- (admit);
\draw[arr] (admit) -- node[right,font=\footnotesize] {yes} (compact);
\draw[arr] (compact) -- (trim);
\draw[arr] (compact.east) -- (publish.west);
\draw[arr] (publish.north) -- (shared.south);
\draw[arr,dashed] (shared.west) -- ++(-.35,0) |- (admit.east);
\node[side,minimum height=1.1cm] (discard) at (6.3,-6)
  {Discard candidate\\Scoring complete};
\draw[arr] (admit.west) -- node[above,font=\footnotesize] {no} ++(-.5,0) |- ([yshift=-.9cm]trim.south)
  -| (discard.south);
\end{tikzpicture}

\caption{Shared complete-key filtering for one scan group. A candidate's
full primary score is computed before admission. A valid $c$-candidate
witness set can tighten the per-query bound; a stale but complete bound
only weakens filtering. Final prefix trimming reduces the valid merge
input, while the allocated workspace capacity remains unchanged.
The whole-key-read premise is qualified in the text.}
\label{fig:5-shared-bound}
\end{figure}

One thread per block reads the query's bound for each chunk. After
compaction, it publishes the local $c$th key through 64-bit
\texttt{atomicMin} when enough valid entries exist. At scan completion,
another publication and a binary search retain only the sorted prefix
no later than the observed bound. There is no cross-block barrier or
wait for other groups. A late update can leave extra entries in another
group's output, which the exact merge safely removes. Bounds are reset
before every search in the same stream, and the merge runs after the
scan kernel finishes.

The read premise deserves a separate qualification. The archived code
publishes atomically but polls through \texttt{volatile} 64-bit loads.
CUDA does not guarantee atomicity or inter-thread synchronisation for
\texttt{volatile} accesses \cite{cuda_programming_guide}. Portable
realisation of the proposition requires device-scope atomic reads as
well. The recorded 48 paired checks in \Cref{sec:7-4} support observed
agreement for the tested build, not a language-level concurrency proof.
An atomic-read revision and its GPU revalidation have not been performed.

\Cref{fig:5-shared-bound} also identifies what this optimisation saves.
Every scanned candidate still loads its primary code and computes its
complete primary score. The bound avoids some buffer appends, local
sorts and intermediate writes, and shortens the lists consumed by the
merge. It does not implement partial-distance early termination or
reduce the number of primary scores evaluated. Its benefit competes
with polling and atomic-update overhead. The observed paired speedup
range includes a slight slowdown, as reported in \Cref{sec:7-4}.

\subsection{Choosing parallelism without changing the budget}

For batch size $b$, grouped scanning launches $bG$ blocks instead of $b$.
This can expose otherwise unavailable parallelism, but duplicates table
staging and local selection and introduces a merge. The fixed-group
design therefore cannot be uniformly preferable. Smaller static quotas
would reduce intermediate storage, but require a certificate or fallback
to preserve coverage; the pilot's concentrated survivor distribution
motivated using witnessed thresholds instead (\Cref{sec:7-4}).

The implemented scheduler uses a fill proxy and a workload-skew
statistic. Let $T$ be scan threads per block, $H$ the number of SMs times
the maximum resident threads per SM, and
$v=\operatorname{sd}(n_c)/\operatorname{mean}(n_c)$ the candidate-count
coefficient of variation. Its rule is
\begin{align}
 G_{\rm fill}
 &=\max\!\left(\{1\}\cup
       \{g\in\{2,4,8\}:bgT\leq0.6H\}\right),
       \label{eq:5-fill-rule}\\
 G_{\rm skew}&=
   \begin{cases}4,&v\geq0.3,\\1,&v<0.3,\end{cases}
 \qquad
 G_{\rm chosen}=\max(G_{\rm fill},G_{\rm skew}).
 \label{eq:5-skew-rule}
\end{align}
The skew rule may override the fill limit. In the evaluation harness,
$v$ is calibrated using 512 database vectors as pseudoqueries and a
fixed routing configuration; it uses no ground-truth labels. This is
distinct from Chapter 6's sampled output-agreement calibration. Once
calibrated, scheduling uses the batch size and stored statistic without
an additional per-search host readback. The constants were fitted on
four pilot configurations and then evaluated on further configurations.

When $G_{\rm chosen}=1$, the implementation selects the faster
query-centric path, labelled $G=0$. It leaves the index, probe count
and refinement budget unchanged, but inherits that path's distance-only
cutoff ties. The common-total-order premise must therefore be checked
before claiming identical survivor sets across all scheduler choices.
\Cref{sec:7-6} reports the separate integrated-output checks and the
scheduler's measured gains and regret.

The rule is a heuristic rather than an occupancy calculation or a
latency oracle: it omits register pressure, shared-memory residency,
budget-dependent compaction cost and merge cost. A frozen pseudoquery
statistic can also misrepresent future queries. In the recorded campaign
the largest regret occurs at batch 256, where a fixed-group alternative
is about $1.217\times$ faster. Thus the design contribution is a tested
combination of parallel candidate scanning, witnessed filtering and
workload-dependent execution, not a new universal top-$c$ identity or a
guarantee of optimal scheduling. Refinement and final top-$k$ consume
the resulting valid survivors as described in \Cref{sec:5-6}.
